% Source: https://artofproblemsolving.com/community/c3h1887431_geometry
% Stable identity: aops-topic-1887431-q01
\textbf{Power of the point $B$.} The line $BA$ is tangent to the circle at $E$, and the line $BC$ meets the circle at $D$ and $C$. So
\[
BE^2=BD\cdot BC=8\cdot11=88 .
\]
\textbf{A lemma.} If $P$ moves along a line $\ell$ perpendicular to $DE$, then $PE^2-PD^2$ stays constant. Indeed, take coordinates with $DE$ on the $x$-axis, $D=(d,0)$, $E=(e,0)$, and $\ell$ the line $x=k$. For $P=(k,y)$,
\[
PE^2-PD^2=(k-e)^2+y^2-(k-d)^2-y^2=(k-e)^2-(k-d)^2,
\]
which does not depend on $y$.

\textbf{Conclusion.} Both $B$ and $P$ lie on the line through $B$ perpendicular to $DE$, so
\[
PE^2-PD^2=BE^2-BD^2=88-64=24 .
\]
With $PE=7$ this gives $PD^2=49-24=25$, so $DP=\boxed{5}$.

\begin{asy}
usepackage("amsmath");
size(6.5cm);
defaultpen(fontsize(10pt));
pair B=(0,0),C=(11,0),D=(8,0);
pair Ee=sqrt(88)*dir(50);
pair A=B+1.3*(Ee-B);
pair u=unit(Ee-D), w=(-u.y,u.x);
pair X=D+dot(B-D,u)*u;
real t2=dot(w,Ee)-sqrt(dot(w,Ee)^2-39);
pair P=t2*w;
pair O=(9.5,(dot(Ee,Ee)-19*Ee.x+88)/(2*Ee.y));
draw(A--B--C--cycle,linewidth(0.9));
draw(circle(O,abs(D-O)),gray(0.55));
draw(D--Ee,linewidth(0.8));
draw((B-0.06*(P-B))--(P+0.06*(P-B)),dashed+gray(0.3));
draw(P--Ee,blue+linewidth(0.9));draw(P--D,red+linewidth(1));
real s=0.4; draw((X+s*u)--(X+s*u-s*unit(X-B))--(X-s*unit(X-B)),linewidth(0.5));
dot(B);dot(C);dot(D);dot(Ee);dot(P);dot(A);
label("$A$",A,N);label("$B$",B,SW);label("$C$",C,SE);label("$D$",D,S);label("$E$",Ee,NW);label("$P$",P,S);
label("$7$",Ee+0.75*(P-Ee),NE,blue);label("$5$",(P+D)/2,NW,red);
label("$8$",(B+D)/2+(0,-0.05),S);label("$3$",(D+C)/2,S);
label("\mbox{}",(min(currentpicture,user=true)+max(currentpicture,user=true))/2,black); // keeps the text after the figure black when LaTeX typesets the labels inline
\end{asy}

(In the figure the dashed line is the perpendicular to $DE$ through $B$; the second point of this line at distance $7$ from $E$ also gives $DP=5$.)
